%=====================================================
\section{Inclusion-Exclusion Principle}
%=====================================================

%-----------------------------------------------------
\begin{frame}{Inclusion-Exclusion Principle}

\textbf{Definition}

\vspace{0.3cm}

The Inclusion-Exclusion Principle is used to determine the
cardinality of the union of overlapping sets without counting
common elements more than once.

\vspace{0.5cm}

For two finite sets,

\[
|A\cup B|
=
|A|
+
|B|
-
|A\cap B|
\]

\end{frame}

%-----------------------------------------------------

\begin{frame}{Example:}

Suppose an autonomous driving model detects
\[
A=
\{\text{Pixels predicted as Road}\}
\]
\[
B=
\{\text{Ground Truth Road Pixels}\}
\]

Suppose
\[
|A|=1200
\]
\[
|B|=1100
\]
\[
|A\cap B|=950
\]
Then
\[
|A\cup B|
=
1200+1100-950
=
1350
\]

The union represents all pixels that are
either predicted as road or actually belong
to the road region.

\end{frame}

%-----------------------------------------------------
\begin{frame}{General Inclusion-Exclusion Formula}

For \(n\) finite sets,

\vspace{0.3cm}

\[
\left|
\bigcup_{i=1}^{n}
A_i
\right|
=
\sum |A_i|
-
\sum |A_i\cap A_j|
+
\sum |A_i\cap A_j\cap A_k|
-\cdots
+(-1)^{n+1}
\left|
\bigcap_{i=1}^{n}
A_i
\right|
\]

\end{frame}